\documentclass[11pt]{article}
\usepackage[margin=1in]{geometry}
\usepackage{natbib}
\usepackage{booktabs}
\usepackage{amsmath}
\usepackage{graphicx}
\usepackage{hyperref}
\usepackage{authblk}

\title{Noncanonical Inflammasome Assembly Requires Caspase-11 Catalytic Activity and Intra-Molecular Autoprocessing}

\author[1]{Authors}
\affil[1]{Division of Immunology and Molecular Medicine, University of California, Berkeley, Berkeley, CA, USA}

\date{}

\begin{document}
\maketitle

\begin{abstract}
Inflammatory caspases are cysteine protease zymogens whose activation following infection or cellular damage occurs within supramolecular organizing centers (SMOCs) known as inflammasomes. Inflammasomes are large oligomeric complexes that recruit caspases to undergo proximity-induced autoprocessing, leading to formation of the enzymatically active form that cleaves downstream targets. Binding of bacterial LPS to its cytosolic sensor, caspase-11 (Casp11), promotes Casp11 aggregation within a high molecular weight complex known as the noncanonical inflammasome, where it is activated to cleave gasdermin~D and induce pyroptosis. However, the cellular correlates of Casp11 oligomerization and whether Casp11 forms an LPS-induced SMOC within cells remain unknown. Using fluorescently labeled Casp11, we found that LPS transfection of macrophages induced Casp11 speck formation, providing direct evidence that Casp11 forms LPS-induced specks in macrophages. Unexpectedly, we found that catalytic activity was required for Casp11 to form LPS-induced specks in macrophages. Importantly, both catalytic activity and autoprocessing were required for Casp11 speck formation in an ectopic expression system, and Casp11 processing via an exogenous protease was sufficient to induce Casp11 speck formation. These data reveal a previously undescribed role for Casp11 catalytic activity and self-cleavage in noncanonical inflammasome assembly and indicate that Casp11 catalytic activity and autoprocessing occur upstream of, and mediate, Casp11 oligomerization. Our data provide new insight into the molecular requirements for assembly of Casp11 noncanonical inflammasome complexes.
\end{abstract}

\section{Introduction}

The mammalian innate immune system relies on evolutionarily conserved pattern recognition receptors (PRRs) to detect microbe-derived pathogen-associated molecular patterns (PAMPs) and rapidly initiate antimicrobial responses \citep{janeway1989,medzhitov2002}. Such rapid responses are critical for efficient host defense in the face of pathogens that have the capacity to replicate rapidly and suppress or evade immune detection. A common feature of PRR activation is their inducible assembly into higher-order oligomeric protein complexes, termed supramolecular organizing centers (SMOCs), which mediate the effector function of these pathways \citep{kagan2014,broz2016}.

Inflammatory caspases are zymogens that contain a C-terminal enzymatic domain and an N-terminal caspase activation and recruitment domain (CARD), which promotes SMOC formation through homotypic interactions \citep{martinon2009,rathinam2012}. The C-terminal enzymatic domain is comprised of large (17--20 kDa) and small (10--12 kDa) subunits separated by an interdomain linker (IDL), which contains the self-cleavage site that undergoes proteolytic processing during caspase activation \citep{boatright2003}.

Caspase-11 (Casp11) serves as both the sensor and effector of the noncanonical inflammasome pathway. Binding of cytosolic LPS to Casp11 promotes its oligomerization into high molecular weight complexes \citep{shi2014,hagar2013,kayagaki2011}. Activated Casp11 cleaves gasdermin~D (GSDMD) to induce pyroptotic cell death \citep{kayagaki2015,shi2015}. However, the cellular correlates of Casp11 oligomerization and whether Casp11 forms an LPS-induced SMOC within cells have remained unknown. Previous work has established that Casp11 auto-proteolysis is important for noncanonical inflammasome activation \citep{lee2018,ross2018}, but the relationship between catalytic activity, autoprocessing, and higher-order complex assembly has not been defined.

Here, we use fluorescently tagged Casp11 constructs to directly visualize inflammasome assembly and dissect the molecular requirements for Casp11 oligomerization in both primary macrophages and HEK293T cells.

\section{Results}

\subsection{Caspase-11 catalytic activity is required for cytosolic LPS-induced speck formation}

Upon sensing cytosolic LPS, Casp11 oligomerizes into high molecular weight complexes known as noncanonical inflammasomes \citep{shi2014}. To better understand the oligomerization dynamics of the Casp11 inflammasome, we generated lentiviral constructs encoding C-terminal fusions of mCherry with full-length Casp11 (Casp11$^{\text{WT}}$-mCherry) and the corresponding C254A mutant lacking catalytic activity, and transduced these into Casp11$^{-/-}$ bone marrow-derived macrophages (BMDMs).

Casp11$^{\text{WT}}$-mCherry but not Casp11$^{\text{C254A}}$-mCherry restored dose-dependent GSDMD cleavage and cytotoxicity in Casp11$^{-/-}$ BMDMs (Figure~1B,C). Confocal microscopy revealed that both constructs exhibited diffuse cytoplasmic staining in the absence of LPS stimulation (Figure~1D). Importantly, Casp11$^{\text{WT}}$-mCherry coalesced into large mCherry specks in response to LPS transfection, consistent with the concept that Casp11 assembles into an LPS-induced higher-order complex, or SMOC (Figure~1D). Casp11$^{\text{WT}}$-mCherry exhibited dose-dependent speck formation paralleling dose-dependent toxicity (Figure~1E). Surprisingly, Casp11$^{\text{C254A}}$-mCherry did not form specks in response to LPS, even at higher doses, despite being expressed at similar or slightly higher levels (Figure~1D,E). These findings imply that CARD-dependent binding of Casp11 to LPS is not sufficient to mediate speck formation and that catalytic activity facilitates assembly of the higher-order noncanonical inflammasome.

\subsection{Caspase-11 catalytic activity is required for spontaneous oligomerization in HEK293T cells}

To further dissect these mechanisms, we employed HEK293T cells, an extensively used system for defining the molecular basis of inflammasome assembly \citep{kayagaki2015,rauch2017,shi2014,tenthorey2014}. Casp11$^{\text{WT}}$-mCherry spontaneously assembled into large aggregates reminiscent of ASC-containing specks \citep{stutz2013,tzeng2016}, whereas Casp11$^{\text{C254A}}$-mCherry remained diffuse within the cytosol (Figure~2A,B).

The pan-caspase inhibitor Z-VAD-FMK (zVAD) blocked spontaneous speck assembly and autoprocessing by Casp11$^{\text{WT}}$-mCherry in a dose-dependent manner (Figure~2C,D). Consistently, Casp11$^{\text{C254A}}$-mCherry did not aggregate into visible specks and its distribution was not altered by zVAD. Importantly, catalytic activity was not required for Casp11 self-association per se, as Casp11$^{\text{WT}}$ and Casp11$^{\text{C254A}}$ associate both with themselves and each other in reciprocal co-IP studies (Extended Data Figure~2C). Speck formation also occurred with a Casp11-citrine fusion but not with a CARD-only-citrine fusion, providing further evidence that the enzymatic domain is important in Casp11 oligomerization (Extended Data Figure~2D).

\subsection{Caspase-11 processing at the interdomain linker is necessary but not sufficient for oligomerization}

We hypothesized that Casp11 autoprocessing might be important in its oligomerization. Casp11$^{\text{D285A}}$-mCherry, which contained a mutation in the IDL autoprocessing site, displayed impaired GSDMD cleavage and cell death \citep{lee2018,ross2018}, and significantly reduced speck formation compared to Casp11$^{\text{WT}}$-mCherry (Figure~3A,B).

To test whether autoprocessing at the IDL is sufficient to induce oligomerization, we replaced the endogenous Casp11 cleavage sequence with the TEV protease consensus sequence (ENLYFQ/G) to generate a TEV-cleavable Casp11 construct (Casp11$^{\text{TEV}}$-mCherry; Figure~3C). Co-expression of TEV protease with Casp11$^{\text{TEV}}$-mCherry resulted in dose-dependent cleavage and induced speck formation, demonstrating that inducible TEV-dependent cleavage could rescue oligomerization of the autoprocessing-deficient mutant (Figure~3D--F). However, TEV protease-induced processing failed to rescue speck formation in C254A Casp11$^{\text{TEV}}$-mCherry, suggesting that Casp11 catalytic activity also functions at an alternate site or that catalytic residue Cys-254 plays an additional role in inflammasome assembly.

\subsection{Wild-type caspase-11 rescues speck formation of catalytically inactive caspase-11 independently of trans-processing}

We co-transfected fixed doses of labeled Casp11$^{\text{C254A}}$-mCherry with increasing amounts of unlabeled wild-type Casp11 (Figure~4A). Expression of unlabeled Casp11 significantly increased speck formation by Casp11$^{\text{C254A}}$-mCherry (Figure~4C,D). Unexpectedly, however, co-expression of untagged Casp11 did not lead to significant levels of Casp11$^{\text{C254A}}$-mCherry processing except at the highest concentrations (Figure~4B), indicating that Casp11 does not induce significant levels of trans-processing in the course of oligomerization into specks.

Additional mutation of the IDL autoprocessing site of the Casp11$^{\text{C254A}}$-mCherry (C254A/D285A double mutant) did not affect its ability to form specks in the presence of co-transfected WT unlabeled Casp11 (Figure~4E--H). These data demonstrate that both catalytically inactive and inactive-uncleavable Casp11 can be recruited into oligomeric inflammasome complexes in the presence of catalytically active Casp11, which must undergo intramolecular processing to induce formation of a Casp11 SMOC.

\subsection{Caspase-11 \textit{cis} autoprocessing mediates noncanonical inflammasome assembly}

To directly test whether rescue is driven by \textit{cis} autoprocessing, we co-transfected labeled Casp11$^{\text{C254A}}$-mCherry with either catalytically inactive Casp11$^{\text{C254A}}$ or non-cleavable Casp11$^{\text{D285A}}$ unlabeled species (Figure~5A). WT Casp11 could undergo autoprocessing, which was absent in cells transfected with either mutant (Figure~5B). Critically, WT unlabeled Casp11 again rescued speck formation by catalytically inactive Casp11-mCherry, but neither unlabeled Casp11$^{\text{C254A}}$ nor Casp11$^{\text{D285A}}$ could restore speck formation for Casp11$^{\text{C254A}}$-mCherry (Figure~5C,D).

These findings demonstrate that \textit{cis} autoprocessing of Casp11 results in assembly of a functional oligomeric complex to which catalytically deficient reporter mutants can be recruited (Figure~5E). Altogether, these results demonstrate a key role for Casp11 catalytic activity and \textit{cis} autoprocessing in oligomerization and higher-order assembly of the Casp11 inflammasome, and indicate that these activities occur upstream of, rather than as a terminal consequence of, inflammasome assembly.

\section{Discussion}

Our study provides direct visualization of Casp11 inflammasome specks in macrophages and reveals an unexpected requirement for catalytic activity and intramolecular autoprocessing in noncanonical inflammasome assembly. The prevailing model of inflammasome activation posits that caspase oligomerization occurs first, followed by proximity-induced autoprocessing \citep{martinon2009,kagan2014}. Our findings fundamentally revise this model for the noncanonical pathway by demonstrating that Casp11 catalytic activity and autoprocessing at the IDL are prerequisites for, rather than consequences of, higher-order oligomerization.

Several lines of evidence support this conclusion. First, catalytically inactive Casp11 fails to form specks in both macrophages and HEK293T cells despite retaining the ability to self-associate at lower-order levels (co-IP). Second, pharmacological caspase inhibition blocks speck formation. Third, a non-cleavable mutant (D285A) shows impaired speck formation. Fourth, TEV-mediated cleavage of the IDL can rescue speck formation only in the presence of intact catalytic activity. Fifth, WT Casp11 rescues speck formation of catalytically inactive mutants through \textit{cis} (intramolecular) but not \textit{trans} (intermolecular) processing.

We propose a model in which Casp11 undergoes initial low-order interactions with LPS through its CARD domain, followed by intramolecular autoprocessing that generates a conformational change permissive for higher-order SMOC assembly. This is analogous to the requirement for proteolytic processing in NLRP1 and CARD8 inflammasome activation \citep{chavarria2016,sharif2021}, and consistent with the structural paradigm in which IDL cleavage releases constraints that prevent caspase oligomerization \citep{boatright2003,shang2021}.

The finding that the catalytic cysteine (C254) contributes to speck formation through mechanisms independent of IDL cleavage suggests additional roles for this residue. One possibility is that C254 mediates cleavage at an alternative site; another is that the catalytic cysteine participates directly in structural interactions required for higher-order assembly \citep{boucher2018}. Future studies will be needed to distinguish between these possibilities.

\section{Materials and Methods}

\subsection{Cell culture}
HEK293T cells were purchased from ATCC and grown in complete DMEM (supplemented with 10\% v/v FBS, 10~mM HEPES, 10~mM sodium pyruvate, 1\% penicillin/streptomycin) at 37$^\circ$C with 5\% CO$_2$. Murine bone marrow progenitors were harvested from femurs and hip bones of 8--12 week-old C57BL/6 or Casp11$^{-/-}$ mice and differentiated into BMDMs using 30\% L929-conditioned complete DMEM for 7 days \citep{bjanes2021}.

\subsection{Cloning and mutagenesis}
The Casp11 coding sequence was fused to mCherry on its C-terminus on pTwist Lenti-SFFV-WPRE lentiviral vector (Twist Biosciences) or mammalian expression vector pReceiver-M56 (Genecopoeia). Catalytic (C254A) and cleavage (D285A) mutants were generated by site-directed mutagenesis (Q5 SDM Kit, New England BioLabs). TEV-cleavable constructs were generated by engineering the TEV protease consensus cleavage sequence (ENLYFQ/G) immediately upstream of cleavage-deficient (D285A) Casp11-mCherry.

\subsection{Microscopy and speck quantification}
Cells were fixed and stained with Hoechst for nuclear visualization. Imaging was performed by fluorescence microscopy. Speck formation was quantified from 4 image frames (100--150 cells) per well and 3 wells per condition for a total of 300--450 cells per condition. Dose-response curves were plotted using non-linear regression (Log2; three parameters).

\subsection{Western blotting}
Whole cell lysates were harvested 12 hours post-transfection and immunoblotted for mCherry or $\beta$-actin as loading control. GSDMD processing was detected using anti-GSDMD antibody.

\subsection{Statistical analysis}
All error bars represent mean $\pm$ SEM of triplicate wells; data are representative of 2--3 independent experiments. Data were analyzed by two-way ANOVA with Sidak's multiple comparison test. Significance thresholds: $*p < 0.05$, $**p < 0.01$, $***p < 0.001$, $****p < 0.0001$.

\bibliographystyle{plainnat}
\bibliography{references}

\end{document}
